\section{Introduction}\label{sec:intro}
The Weil criterion expresses the Riemann hypothesis as positivity of an
explicit quadratic form on compactly supported test functions. An inverse
bulk approach asks for a quantum system, a boundary or defect preparation,
and an independently positive physical pairing whose pullback is exactly
that form. Interactions are compatible with this objective: a preparation
linear in the input produces a quadratic state norm even when the
underlying dynamics is interacting. What must be derived is the entire
pairing, with its physical adjoint and normalization.

Here ``boundary'' means the locus or observable sector through which
the arithmetic input enters the system. It may be a spatial boundary,
a line defect, an interface, a gluing cut, or a module for a protected
operator algebra. This flexibility permits the use of sphere and Schur
quantization, where explicit representations and positive pairings are
available \cite{Sphere,Schur,RG}. Conformal symmetry can constrain a chosen
defect, but does not identify arbitrary deformations or source maps.
Pure four-dimensional Yang--Mills is also a possible long-term construction
target; no identification with it is made in this paper.

\clearpage
\subsection{Principal result and its limitation}
The most concrete arithmetic result concerns the elementary Schur
representation. For a fixed quantization parameter $0<q<1$, we construct
an electric dressing and one magnetic insertion whose norm on Laurent
polynomials $h$ is
\begin{equation}\label{eq:intro-norm}
 \norm{\Phi_{q;r,d}[h]}^2
 =\int_{S^1}|h(z)|^2
 \frac{dr(1+r)}{1-r}\frac{|1-z|^2}{|1-rz|^2}\frac{\dd\theta}{2\pi},
 \qquad z=e^{i\theta},\quad 0<r<1,
\end{equation}
with $d>0$. Its Fourier coefficients are
\begin{equation}\label{eq:intro-coeff}
 \frac{dr(1+r)}{1-r}\frac{|1-z|^2}{|1-rz|^2}
 =\frac{2dr}{1-r}-d\sum_{n\geq1}r^n(z^n+z^{-n}).
\end{equation}
Thus $r=p^{-1/2}$ and $d=\log p$ give every required single-prime
repetition coefficient. The arithmetic parameters enter the preparation;
$q$ is held fixed. Theorem~\ref{thm:schur-prime} proves this identity and
the domain statement needed for its magnetic insertion. An explicit
unitary decomposition then relates the circle calculation to the real
logarithmic coordinate of the Weil form.

The positive contact $2dr/(1-r)$ is part of the result. This is a Schur
realization of the positive prime reference already present in the
project's conservative return construction. Neither the reference form
nor the geometric Fourier identity is claimed as new. The additional
content is the specified operator preparation, its source normalization,
its graph-domain justification, and its compatibility with one fixed
Schur parameter. A physical interface implementing the continuous-input
lift has not been constructed.

\subsection{Other results and the remaining question}
The Gaussian real boundary module identifies the archimedean
quarter-shift and oscillator levels, but its natural state-smearing norm
has a different spectral weight. In the interacting $T[SU(2)]$ sphere
model we evaluate an infinite family of charged-state norms and show
that bare powers do not have geometric return ratios. A natural electric
preparation in the conformal $SU(2)$ theory with four hypermultiplets
also fails, as does the undressed elementary Schur return prescription.
These are exclusions of specified preparations, not of the theories.

Two further results govern the new dressed source. Its finite magnetic
domain does not permit unrestricted use of higher-power identities.
Moreover, gluing the prime filters through fixed output vectors produces
an atom at $\log(3/2)$ unless those vectors are orthogonal
(Proposition~\ref{prop:mixed}). Orthogonality restores the independent
positive contacts. A successful joint construction must therefore change
the preparation or the gluing operation in a substantive way.

The background and normalization are developed in
Section~\ref{sec:weil}, independently of any companion manuscript.
Section~\ref{sec:positive} states the conditional construction-to-RH
implication and a general source-regularity restriction.
Sections~\ref{sec:sphere}--\ref{sec:gluing} contain the main calculations.
Section~\ref{sec:outlook} formulates the next complete fitting problem.
The appendices retain positive gamma refinement and earlier gauge and
rational-feedback tests, with proofs sufficient to use them without
consulting the research notes.

The classical Weil criterion, gamma identities, and established
representation formulas are explicitly attributed. The comparison
calculations are working results, not a claim of literature priority.
In particular, an exact identity between local phase derivatives and
the prime terms is not a positivity theorem. No candidate in this draft
matches the complete first-prime Weil form, and no all-support result
follows from the local norm in \eqref{eq:intro-norm}.
